\documentclass[10pt,a4paper]{amsart}
\usepackage{amsmath,amssymb,amsthm,graphicx}
\usepackage{verbatim}
\newenvironment{reference}{\comment}{\endcomment}
%\newenvironment{reference}{}{}
\newenvironment{slogan}{\comment}{\endcomment}
\newenvironment{history}{\comment}{\endcomment}

% For commutative diagrams we use Xy-pic
\usepackage[all]{xy}

% We use 2cell for 2-commutative diagrams.
\xyoption{2cell}
% Theorem environments.
%
\theoremstyle{plain}
\newtheorem{theorem}[subsection]{Theorem}
\newtheorem{proposition}[subsection]{Proposition}
\newtheorem{lemma}[subsection]{Lemma}

\theoremstyle{definition}
\newtheorem{definition}[subsection]{Definition}
\newtheorem{example}[subsection]{Example}
\newtheorem{exercise}[subsection]{Exercise}
\newtheorem{situation}[subsection]{Situation}

\theoremstyle{remark}
\newtheorem{remark}[subsection]{Remark}
\newtheorem{remarks}[subsection]{Remarks}

\numberwithin{equation}{subsection}

% Macros
%
\def\lim{\mathop{\mathrm{lim}}\nolimits}
\def\colim{\mathop{\mathrm{colim}}\nolimits}
\def\Spec{\mathop{\mathrm{Spec}}}
\def\Hom{\mathop{\mathrm{Hom}}\nolimits}
\def\Ext{\mathop{\mathrm{Ext}}\nolimits}
\def\SheafHom{\mathop{\mathcal{H}\!\mathit{om}}\nolimits}
\def\SheafExt{\mathop{\mathcal{E}\!\mathit{xt}}\nolimits}
\def\Sch{\mathit{Sch}}
\def\Mor{\mathop{\mathrm{Mor}}\nolimits}
\def\Ob{\mathop{\mathrm{Ob}}\nolimits}
\def\Sh{\mathop{\mathit{Sh}}\nolimits}
\def\NL{\mathop{N\!L}\nolimits}
\def\CH{\mathop{\mathrm{CH}}\nolimits}
\def\proetale{{pro\text{-}\acute{e}tale}}
\def\etale{{\acute{e}tale}}
\def\QCoh{\mathit{QCoh}}
\def\Ker{\mathop{\mathrm{Ker}}}
\def\Im{\mathop{\mathrm{Im}}}
\def\Coker{\mathop{\mathrm{Coker}}}
\def\Coim{\mathop{\mathrm{Coim}}}

% Boxtimes
%
\DeclareMathSymbol{\boxtimes}{\mathbin}{AMSa}{"02}

%
% Macros for moduli stacks/spaces
%
\def\QCohstack{\mathcal{QC}\!\mathit{oh}}
\def\Cohstack{\mathcal{C}\!\mathit{oh}}
\def\Spacesstack{\mathcal{S}\!\mathit{paces}}
\def\Quotfunctor{\mathrm{Quot}}
\def\Hilbfunctor{\mathrm{Hilb}}
\def\Curvesstack{\mathcal{C}\!\mathit{urves}}
\def\Polarizedstack{\mathcal{P}\!\mathit{olarized}}
\def\Complexesstack{\mathcal{C}\!\mathit{omplexes}}
% \Pic is the operator that assigns to X its picard group, usage \Pic(X)
% \Picardstack_{X/B} denotes the Picard stack of X over B
% \Picardfunctor_{X/B} denotes the Picard functor of X over B
\def\Pic{\mathop{\mathrm{Pic}}\nolimits}
\def\Picardstack{\mathcal{P}\!\mathit{ic}}
\def\Picardfunctor{\mathrm{Pic}}
\def\Deformationcategory{\mathcal{D}\!\mathit{ef}}
\usepackage{cleveref}
\crefname{originalequation}{original equation}{original equations}
\Crefname{originalequation}{Original equation}{Original equations}
\makeatletter
\newlabel{equation-FM-QCoh}{{0FZC}{source}{}{sourcepointer-1}{}}
\newlabel{equation-FM-QCoh@cref}{{[originalequation][0][]0FZC}{[1][0][]source}{}{}{}}
\makeatother
\crefname{originalreference}{original reference}{original references}
\Crefname{originalreference}{Original reference}{Original references}
\makeatletter
\newlabel{example-functor-quasi-coherent}{{0FZB}{source}{}{sourcepointer-2}{}}
\newlabel{example-functor-quasi-coherent@cref}{{[originalreference][0][]0FZB}{[1][0][]source}{}{}{}}
\makeatother
\crefname{originalreference}{original reference}{original references}
\Crefname{originalreference}{Original reference}{Original references}
\makeatletter
\newlabel{morphisms-lemma-base-change-affine}{{01SD}{source}{}{sourcepointer-3}{}}
\newlabel{morphisms-lemma-base-change-affine@cref}{{[originalreference][0][]01SD}{[1][0][]source}{}{}{}}
\makeatother
\crefname{originalreference}{original reference}{original references}
\Crefname{originalreference}{Original reference}{Original references}
\makeatletter
\newlabel{coherent-lemma-relative-affine-vanishing}{{01XC}{source}{}{sourcepointer-4}{}}
\newlabel{coherent-lemma-relative-affine-vanishing@cref}{{[originalreference][0][]01XC}{[1][0][]source}{}{}{}}
\makeatother
\crefname{originalreference}{original reference}{original references}
\Crefname{originalreference}{Original reference}{Original references}
\makeatletter
\newlabel{morphisms-lemma-affine-equivalence-modules}{{01SB}{source}{}{sourcepointer-5}{}}
\newlabel{morphisms-lemma-affine-equivalence-modules@cref}{{[originalreference][0][]01SB}{[1][0][]source}{}{}{}}
\makeatother
\crefname{originalreference}{original reference}{original references}
\Crefname{originalreference}{Original reference}{Original references}
\makeatletter
\newlabel{schemes-lemma-equivalence-quasi-coherent}{{01IB}{source}{}{sourcepointer-6}{}}
\newlabel{schemes-lemma-equivalence-quasi-coherent@cref}{{[originalreference][0][]01IB}{[1][0][]source}{}{}{}}
\makeatother
\crefname{originallemma}{original lemma}{original lemmas}
\Crefname{originallemma}{Original lemma}{Original lemmas}
\makeatletter
\newlabel{lemma-functor-on-modules}{{0GNQ}{source}{}{sourcepointer-7}{}}
\newlabel{lemma-functor-on-modules@cref}{{[originallemma][0][]0GNQ}{[1][0][]source}{}{}{}}
\makeatother
\title{R-linear right exact direct-sum-commuting functors QCoh(X) → QCoh(Y) for affine X are QCoh(X ×\_R Y)}
\author{The Stacks Project Authors}
\date{stacks-project:functors.tex:local-snapshot-20261012; GNU FDL 1.2 or later}
\begin{document}
\maketitle
\noindent\textit{Editorial note.} All delivered mathematical text is the author's own, copied verbatim from the pinned Stacks Project chapter file functors.tex; no statement or proof step was rewritten, reordered or removed.\par\medskip
\noindent\textit{Editorial note.} Span recomputation by label pairing (LANE\_BRIEF 10.2): in this chapter the author's proof environment follows the statement environment as a standalone block; the spec covers the complete environment pair (statement span, proof span) as recorded in the pool candidate.\par\medskip
\noindent\textit{Editorial note.} Occupation scan: combined claim\_key + statement\_anchor + generic-line scan over all registration files (active and withdrawn) found 0 hits for functors.tex (fu-occupied.json, 2026-10-14); consistent with the duty manager's six-check dispatch.\par\medskip
\ifdefined\hypertarget\hypertarget{sourcepointer-1}{}\fi
\noindent\textit{Source reference.} Original equation 0FZC: sources/stacks-chow/tags.txt, source lines 19947--19947. Cross-chapter standard prerequisite cited by the author inside the counted span; it is used as a pointer and its own proof is not part of this item.\par\medskip
\ifdefined\hypertarget\hypertarget{sourcepointer-2}{}\fi
\noindent\textit{Source reference.} Original reference 0FZB: sources/stacks-chow/tags.txt, source lines 19946--19946. Cross-chapter standard prerequisite cited by the author inside the counted span; it is used as a pointer and its own proof is not part of this item.\par\medskip
\ifdefined\hypertarget\hypertarget{sourcepointer-3}{}\fi
\noindent\textit{Source reference.} Original reference 01SD: sources/stacks-chow/tags.txt, source lines 2334--2334. Cross-chapter standard prerequisite cited by the author inside the counted span; it is used as a pointer and its own proof is not part of this item.\par\medskip
\ifdefined\hypertarget\hypertarget{sourcepointer-4}{}\fi
\noindent\textit{Source reference.} Original reference 01XC: sources/stacks-chow/tags.txt, source lines 2513--2513. Cross-chapter standard prerequisite cited by the author inside the counted span; it is used as a pointer and its own proof is not part of this item.\par\medskip
\ifdefined\hypertarget\hypertarget{sourcepointer-5}{}\fi
\noindent\textit{Source reference.} Original reference 01SB: sources/stacks-chow/tags.txt, source lines 2332--2332. Cross-chapter standard prerequisite cited by the author inside the counted span; it is used as a pointer and its own proof is not part of this item.\par\medskip
\ifdefined\hypertarget\hypertarget{sourcepointer-6}{}\fi
\noindent\textit{Source reference.} Original reference 01IB: sources/stacks-chow/tags.txt, source lines 1972--1972. Cross-chapter standard prerequisite cited by the author inside the counted span; it is used as a pointer and its own proof is not part of this item.\par\medskip
\ifdefined\hypertarget\hypertarget{sourcepointer-7}{}\fi
\noindent\textit{Source reference.} Original lemma 0GNQ: sources/stacks-chow/tags.txt, source lines 20800--20800. Cross-chapter standard prerequisite cited by the author inside the counted span; it is used as a pointer and its own proof is not part of this item.\par\medskip
\begin{lemma}
\label{lemma-functor-quasi-coherent-from-affine}
Let $R$ be a ring. Let $X$ and $Y$ be schemes over $R$ with $X$ affine.
There is an equivalence of categories between
\begin{enumerate}
\item the category of $R$-linear functors
$F : \QCoh(\mathcal{O}_X) \to \QCoh(\mathcal{O}_Y)$
which are right exact and commute with arbitrary direct sums, and
\item the category $\QCoh(\mathcal{O}_{X \times_R Y})$
\end{enumerate}
given by sending $\mathcal{K}$ to the functor $F$ in (\ref{equation-FM-QCoh}).
\end{lemma}

\begin{proof}
Let $\mathcal{K}$ be an object of $\QCoh(\mathcal{O}_{X \times_R Y})$
and $F_\mathcal{K}$ the functor (\ref{equation-FM-QCoh}). By the discussion in
Example \ref{example-functor-quasi-coherent} we already know that
$F$ is $R$-linear and commutes with arbitrary direct sums.
Since $\text{pr}_2 : X \times_R Y \to Y$ is affine
(Morphisms, Lemma \ref{morphisms-lemma-base-change-affine}) the functor
$\text{pr}_{2, *}$ is exact, see Cohomology of Schemes, Lemma
\ref{coherent-lemma-relative-affine-vanishing}.
Hence $F$ is right exact as well, in other words $F$ is as in (1).

\medskip\noindent
Let $F$ be as in (1). Say $X = \Spec(A)$. Consider the quasi-coherent
$\mathcal{O}_Y$-module $\mathcal{G} = F(\mathcal{O}_X)$.
The functor $F$ induces an $R$-linear map
$A \to \text{End}_{\mathcal{O}_Y}(\mathcal{G})$,
$a \mapsto F(a \cdot \text{id})$. Thus $\mathcal{G}$ is a sheaf of modules over
$$
A \otimes_R \mathcal{O}_Y = \text{pr}_{2, *}\mathcal{O}_{X \times_R Y}
$$
By Morphisms, Lemma \ref{morphisms-lemma-affine-equivalence-modules}
we find that there is a unique quasi-coherent module $\mathcal{K}$
on $X \times_R Y$ such that $F(\mathcal{O}_X) = \mathcal{G} =
\text{pr}_{2, *}\mathcal{K}$ compatible with action
of $A$ and $\mathcal{O}_Y$. Denote $F_\mathcal{K}$ the functor
given by (\ref{equation-FM-QCoh}). There is an equivalence
$\text{Mod}_A \to \QCoh(\mathcal{O}_X)$ sending $A$ to $\mathcal{O}_X$, see
Schemes, Lemma \ref{schemes-lemma-equivalence-quasi-coherent}.
Hence we find an isomorphism $F \cong F_\mathcal{K}$ by
Lemma \ref{lemma-functor-on-modules} because we have an isomorphism
$F(\mathcal{O}_X) \cong F_\mathcal{K}(\mathcal{O}_X)$ compatible with
$A$-action by construction.

\medskip\noindent
This shows that the functor sending $\mathcal{K}$ to $F_\mathcal{K}$
is essentially surjective. We omit the verification of fully faithfulness.
\end{proof}

\end{document}
